\cotefile{93}
\scanpage{165}
\[
 H^*_{\mathrm{dR}}(X/T)\Longleftarrow E_2^{pq}
 =H^p(S,\alpha\longmapsto\mathscr H^q(\Omega^\alpha_{S/T}
 \overset{L}{\otimes}\mathbb R f_*(\Omega^\cdot_{X/S}))),
 \tag*{(12)}
\]
où l'opérateur diff. pour le complexe $\alpha\longmapsto\ldots$ est celui déduit de la connexion de Gauss-Manin de $\mathbb R f_*(\Omega^\cdot_{X/S})$. Dans le cas favorable $\Omega^1_{S/T}$ plat, \inserted{(il suffit $\Omega^1_{S/T}$ plat, par D.~Lazard)} \deleted{\unsure{la flèche précédente}} la suite spectrale devient
\[
\begin{aligned}
 H^*_{\mathrm{DR}}(X/T)\Longleftarrow E_2^{pq}
 &=H^p(S,\alpha\longmapsto\Omega^\alpha_{S/T}\otimes\mathbb R^qf_*(\Omega^\cdot_{X/S}))\\
 &=H^p_{\mathrm{DR}}(S/T,\mathbb R^q f_*(\Omega^\cdot_{X/S})),
\end{aligned}
\tag*{(12')}
\]
où l'expression $H^p_{\mathrm{DR}}$ a un sens, grâce à la connexion de Gauss-Manin de $\mathbb R^qf_*(\Omega^\cdot_{X/S})$.

Comme de juste, \unsure{les suites spectrales|la suite spectrale} (12) \deleted{et (12')} (donc (12')) s'appelle \uline{suite spectrale de Leray} \deleted{\unsure{pour}} \unsure{de} $f$, en cohomologie de De Rham.

\deleted{\unsure{Voici une}} \uline{Application}\quad Soit $n$ un entier, et soit
\[
 \xi\in H^n_{\mathrm{DR}}(X/T)
 \tag*{(13)}
\]
telle que son image dans $E_2^{0,n}=\deleted{\unsure{H|\mathbb H^0}}\operatorname{Ker}(H^0(S,\mathbb R^n f_*(\Omega^\cdot_{X/S}))\deleted{\to H^0}\inserted{\to H^0(S,\mathscr H^n(\Omega^1_{X/S}\overset{L}{\otimes}\mathbb R f_*(\Omega^\cdot_{X/S})))})$ \unsure{par}
\[
 H^n_{\mathrm{DR}}(X/T)\longrightarrow H^0(S,\underbrace{\mathbb R^n f_*(\Omega^\cdot_{X/S})}_{\mathscr H^n_{\mathrm{DR}}(X/S)}),
 \quad\deleted{\unsure{(\to H^0(S,\mathscr H^n\mathbb R f_*(\Omega^\cdot_{X/S}))) }}
 \tag*{(14}
\]
est \uline{nulle} [\uline{NB} cette image est en tous cas tjrs \uline{horizontale}, i.e. son image dans $\Gamma(S,\mathscr H^n(\Omega^1_{X/S}\overset{L}{\otimes}\mathbb R f_*(\Omega^\cdot_{X/S})))$ est nulle). 
% archive sheet number: 164

